% Ajustes de LaTeX para la guía técnica (docs/pdf/build.sh). Se compila con pdfLaTeX.
% Símbolos Unicode que usan los Markdown de docs/ (pdfLaTeX no los trae; las cajas del árbol de carpetas las pone
% pmboxdraw). Si un documento nuevo usa otro símbolo, se agrega aquí.
\usepackage{graphicx} % los diagramas entran como LaTeX directo (guia.lua)
\usepackage{tikz}     % escudos y líneas de la portada (portada.tex)
\usepackage{pmboxdraw}
\usepackage{pifont}
\usepackage{amssymb}
\usepackage{newunicodechar}
\newunicodechar{✓}{\ding{51}}
\newunicodechar{✗}{\ding{55}}
\newunicodechar{↻}{\ensuremath{\circlearrowright}}
\newunicodechar{→}{\ensuremath{\rightarrow}}
\newunicodechar{≥}{\ensuremath{\geq}}
\newunicodechar{≤}{\ensuremath{\leq}}
\newunicodechar{≈}{\ensuremath{\approx}}
\newunicodechar{−}{\ensuremath{-}}
% Bloques de código: cortar las líneas largas en lugar de salirse del margen.
\usepackage{fvextra}
\DefineVerbatimEnvironment{Highlighting}{Verbatim}{breaklines,breakanywhere,commandchars=\\\{\}}
% Tablas: letra un poco menor y celdas con más aire.
\AtBeginEnvironment{longtable}{\small}
\renewcommand{\arraystretch}{1.25}
% Párrafos separados por espacio, sin sangría (como en la vista de GitHub).
\setlength{\parindent}{0pt}
\setlength{\parskip}{0.55em plus 0.1em}
% Encabezado con el capítulo en curso.
\usepackage{fancyhdr}
\pagestyle{fancy}
\fancyhf{}
\fancyhead[L]{\small\nouppercase{\leftmark}}
\fancyhead[R]{\small Guía técnica · TT 2026-B033}
\fancyfoot[C]{\thepage}
\renewcommand{\headrulewidth}{0.4pt}
